\section{Major mechanisms}
\label{sec:mech}
The campaign supports four mechanisms, each backed by matched comparisons and, where stochastic, by seed replicates.

\paragraph{1. Net section times ligament-size-dependent net-section strength.} The strength of a porous sheet is the minimum load-bearing solid fraction across the loading direction multiplied by a net-section strength that depends on the width and alignment of the load-bearing ligaments but not on pore shape: $\approx$30--32\,N/m (80\% of the intrinsic strength) for straight, load-parallel ligaments wider than $\sim$20\,\AA, 17--19\,N/m for round-pore ligaments of 10--13\,\AA, and 12--20\,N/m for random networks. Evidence: porosity independence of the parallel-slit strength, the linear net-section scaling of the square coarse mesh, the equality of square and round coarse pores at equal net section, the master plot of net-section strength versus ligament width, and the holdout predictions (Sec.~\ref{sec:holdout}). Mechanistically, the fine ligaments are weaker because a larger fraction of their bonds is perturbed by the two-coordinated pore edges (edge reconstruction and bond-order changes), an atomic-scale effect that has no continuum analogue at fixed geometry ratio.

\paragraph{2. Load-path orientation.} Ligaments or pores aligned with the load carry the stress without concentration; features across the load reduce the net section and concentrate stress at their ends. The same porosity yields 25.4 vs 4.0\,N/m for slits parallel vs perpendicular to the load, 19.0 vs 8.6\,N/m for ellipses, and 16.1 vs 12.0\,N/m for veins along vs across the load in nested meshes. This is the largest single lever in the design space and acts identically in single-scale and hierarchical structures.

\paragraph{3. Fibre-bundle load transfer sets the fracture mode.} After the first local failure, the released load is shared by the remaining parallel ligaments; failure is progressive only when the share transferred to each survivor stays below its strength margin, i.e. when there are many parallel paths of moderate strength dispersion. Few wide strips (three per cell) cascade even with 15\% dispersion (Stage~4), single flaws and clustered defects run in one event, whereas 8--10 rows of fine ligaments, jittered lattices, distributed vacancies and load-parallel-vein meshes fail row by row and retain 20--60\% of the peak load. Spatial localisation of the damage is a separate property: periodic single-avalanche failures are spatially delocalised, crack-like failures localised.

\paragraph{4. Lattice trapping of cracks.} Beyond $\approx$20\,\AA{} the strength of a cracked or holed sheet becomes independent of flaw size (18--19.5\,N/m for 20--50\,\AA{} flaws in 100--150\,\AA{} cells), because crack advance in the discrete lattice proceeds by trapped bond-breaking events with a barrier; stable bond-by-bond growth precedes the instability. Griffith scaling from the 10 and 20\,\AA{} cracks would have predicted 12--14\,N/m for the 40 and 50\,\AA{} cracks.

\paragraph{What hierarchy does.} A nested mesh behaves as the parallel combination of its finest level (which fails first) and its load-parallel coarse ligaments; at matched porosity it is never stronger than the coarse level alone, and never tougher than the fine level alone, unless the fine level is itself aligned with the load. Scale ratio, hierarchy depth and perpendicular veins are irrelevant or detrimental. Hierarchy is therefore not a mechanism in this model; alignment and ligament width are.

\section{Rejected hypotheses and counter-examples}
\begin{itemize}
\item \textbf{H-A (pore-shape knockdown) -- rejected in its original form.} Circular and square coarse pores at equal net section are equally strong (16.1 vs 17.0\,N/m; predicted 11.3). The net-section factor survived, the shape factor did not.
\item \textbf{H-C (load-path equivalence controls the mode) -- rejected.} Dispersing ligament widths by 15\% lowered strength in both coarse meshes and left the single-avalanche mode unchanged (square) or shortened the tail (round). Replaced by the fibre-bundle load-transfer criterion.
\item \textbf{H-D (disorder optimum for energy absorption) -- not supported after replication.} The Voronoi ``optimum'' at 60\% regularity was a single-seed outlier (14.4 vs 9.2/9.9\,N/m); jittered meshes lose strength without a reproducible gain in work ($1.07\pm0.31$ vs $1.16\pm0.30$\,J/m$^2$ for the ordered lattice across registries).
\item \textbf{H-F$'$ (veins as crack arrestors) -- rejected.} Load-perpendicular veins neither arrest cracks nor add strength (12.0\,N/m, equal to the fine mesh); cracks channel along the columns next to them.
\item \textbf{Flaw halo (P11) -- no shielding.} Pore rings around a hole cost strength relative to the bare hole (16.2 vs 19.1\,N/m); they only reduce the penalty relative to uniformly distributed porosity.
\item \textbf{Gradients (P10) -- detrimental.} Pore-size gradients along or radial to the load are weaker than the uniform mesh (10.1/10.0 vs 11.7\,N/m) because the largest pores define the weakest section.
\item \textbf{Node connectivity (P9) -- minor.} Six-connected triangulated struts are not stronger than four-connected ones; a continuous strut family along the load is what matters (11.5 vs 7.4 vs 6.0\,N/m).
\item \textbf{Unstable and strongly reconstructed candidates.} Every generated geometry relaxed to a mechanically continuous structure, but 6\,\AA{} struts, 20\,\AA{}-vein meshes with residual 2-atom pores, and the 0.40-porosity mesh reconstructed by more than 1\,\AA{} (flagged ``strongly reconstructed''); their results are reported with that flag. No structure fragmented, and none is claimed to be synthesizable.
\item \textbf{Counter-examples worth keeping.} The 45$^\circ$ slit lattice sheds load at 17\% strain and then re-stiffens to $>15$\,N/m at 30\% by ligament rotation (geometric compliance); the 30\,\AA{} crack is not weaker than the 20\,\AA{} crack; the fine period-12 mesh is stronger than the period-16 mesh at equal porosity, but the difference (13.3 vs 11.7\,N/m) is within the $\approx$12\% registry effect measured in Stage~6.
\end{itemize}

\section{Limitations}
\begin{itemize}
\item \textbf{Model, not experiment.} Every number in this report is a property of the screened second-generation REBO potential under athermal quasi-static loading. \rebo{} underestimates the in-plane stiffness of graphene by $\sim$30\% and overestimates its Poisson ratio; intrinsic strengths and the qualitative bond-breaking chemistry are reasonable, but absolute strengths and failure strains should be read as model quantities. No structure is claimed to be synthesizable; generated structures are not new allotropes.
\item \textbf{Athermal, quasi-static, 2D-planar.} Thermal activation, rate effects and out-of-plane deformation (buckling, wrinkling of ligaments) are absent by construction; the sheets are kept planar. The optional MD module allows finite-temperature checks but was not used for the main campaign.
\item \textbf{Periodic supercells.} Loading is applied to periodic cells of 100--120\,\AA; pores and cracks interact with their periodic images, and the coarsest hierarchical level is limited by the cell. Size effects were checked (TEST~15) but not eliminated.
\item \textbf{Uniaxial-stress protocol.} The transverse cell is relaxed to $\sigma_{yy}=0$ by a secant scheme with a trust region; for strongly non-monotonic transverse responses (pristine zigzag graphene above 16\% strain) the followed branch depends on history.
\item \textbf{Bond definition.} Damage is inferred from a geometric criterion ($r<2.0$\,\AA) applied to converged states; events shorter than one strain increment are invisible, and the bisection resolves only the first event to 0.125\%.
\item \textbf{Float32 fast mode.} The campaign used float32 on the Apple GPU; differences to the float64 reference are $10^{-4}$\,eV/\AA{} in forces and $10^{-5}$\,N/m in stress, but MPS reductions are order-non-deterministic at the $10^{-7}$ level, so post-instability branches of a given run are reproducible only statistically (Stage~6 seeds address the structural, not the numerical, variability).
\item \textbf{Descriptors.} Structural descriptors are raster/graph based and coarse (0.5\,\AA{} raster, 8$\times$8 localisation grid); the automatic fracture-mode classification uses fixed thresholds.
\item \textbf{Sampling.} About one hundred simulations cannot cover a twelve-dimensional design space densely; the study identifies mechanisms and regimes, not optima.
\end{itemize}

\section{Threats to validity}
\begin{itemize}
\item \textbf{Circularity.} Hypotheses were formed from Stage~1--2 data and tested on new structures in Stage~4 and on holdouts in Stage~7--8; the prediction files are timestamped and hashed before the corresponding simulations were launched. Retrospective modification of predictions was not performed.
\item \textbf{Confounding.} Comparisons across architectures are controlled for porosity (mass per footprint) and cell size, but not simultaneously for every descriptor (edge-atom fraction, ligament width and net section co-vary with the architecture); the property maps and matched pairs in Stage~4 separate these where possible.
\item \textbf{Numerical artefacts.} Non-converged minimisations are flagged in every record; the effect of tolerance, strain increment and cell size was quantified (TESTS~13--15). Screening switch-on features and the affine cut-off re-stiffening are documented and do not affect relaxed trajectories.
\item \textbf{Selection of top structures.} The showcase set was chosen by predefined categories (reference, strongest, most work, Pareto member, hierarchical, most localised, most progressive, flaw halo, counter-example), not by visual appeal.
\item \textbf{Stochastic variability.} Disordered structures were replicated with $\ge3$ seeds in Stage~6; conclusions about disorder rely on those distributions, individual seeds are reported.
\end{itemize}

\section{Conclusions}
A validated, GPU-resident implementation of screened REBO2 was used to run an autonomous, hypothesis-driven study of how atomic-scale architecture controls the fracture of graphene-derived sheets, with the reference images serving as a source of design principles rather than templates. The findings, all statements about the \rebo{} model under athermal quasi-static uniaxial tension, are:
\begin{enumerate}
\item Strength is governed by the minimum load-bearing solid fraction across the load multiplied by a net-section strength that depends on ligament width and alignment (30--32\,N/m for straight, load-parallel ligaments wider than $\sim$20\,\AA; 17--19\,N/m for 10--13\,\AA{} round-pore ligaments; 12--20\,N/m for random networks), not on pore shape, porosity per se, hierarchy depth or scale ratio.
\item Load-path orientation is the largest lever: identical porosity yields 25 vs 4\,N/m for slits parallel vs perpendicular to the load, and 16 vs 12\,N/m for coarse veins along vs across the load.
\item Hierarchy adds no strength or toughness at matched mass unless its levels are aligned with the load; the only nested designs that beat the coarse-only limit combine load-elongated fine pores with veins (17.5\,N/m with square veins, progressive; 19.4\,N/m with load-parallel veins, stepwise; both 2.0\,J/m$^2$).
\item The fracture mode is a load-transfer (fibre-bundle) property: many parallel ligaments of moderate dispersion fail progressively and retain load, few equivalent wide strips, clustered defects and single flaws fail in one avalanche; strength dispersion alone does not create progressivity.
\item Cracks and holes larger than $\approx$20\,\AA{} sit on a lattice-trapping plateau (18--19.5\,N/m for 20--50\,\AA{} flaws), far above Griffith scaling.
\item Disorder, gradients and flaw halos are costs, not toughening mechanisms, in this model; apparent optima disappeared under seed replication, and the atomic registry of pore edges alone shifts fine-mesh strength by $\approx$12\%.
\item Twelve holdout designs predicted before simulation gave a median strength error of 14\% and the correct single-avalanche/multi-event classification in 10 of 12 cases; the misses identified a chirality-dependent ligament strength, the additivity of alignment effects across scales (best nested design 19.4\,N/m), the weakest-section rule for halos and gradients, and a new mechanism, en-echelon slit coalescence, at which load-path alignment is lost (8.9 instead of the predicted 23.9\,N/m).
\end{enumerate}
The platform, database (\nRuns{} records), structures, trajectories, figures, movies and this report are delivered as one reproducible archive; every simulation can be regenerated from its record.

\section{Damage and connectivity analysis}
\label{sec:damage}
Two carbon atoms are counted as bonded when their distance is below 2.0\,\AA{} (between the equilibrium bond length of 1.42\,\AA{} and the inner cut-off region 1.95--2.25\,\AA{} of \rebo{}; bonds at the tensile instability are 1.7--1.8\,\AA{} long and ruptured pairs relax to more than 2.5\,\AA{} within one increment). Bond sets are compared between consecutive converged states; every lost bond is recorded with the strain, the two atoms and the midpoint of the bond. Persistence is enforced in post-processing (a bond that reappears is not counted as damage). The damage-localisation measure is $L=1-H/H_\mathrm{max}$, where $H=-\sum_i p_i\ln p_i$ is the Shannon entropy of the distribution $p_i$ of damage-event midpoints over an $8\times8$ grid of the periodic cell and $H_\mathrm{max}=\ln 64$; $L=0$ for uniformly distributed damage and $L=1$ if all events fall in one cell. Load-bearing connectivity is tested by percolation of the bond graph across the periodic $x$ boundary (a doubled-cell union-find). Fracture modes are classified automatically from the number of discrete damage steps, the fraction of bonds lost in the largest single step, the post-peak strain range and the post-peak load retention (\texttt{analysis/metrics.py}).

